\section{A literal normalized block}

The smallest convenient equal-mass witness occurs at
\[
 Q=39,\qquad L=7,\qquad h=1092,
 \qquad (p,r)=(67,71).
\]
Both rows have positive primitive multipliers $(1,5,11)$ and therefore six signed
atoms.  Their supports intersect exactly at
\[
 \begin{array}{c|cc}
 \text{residue}&m_p&m_r\\ \hline
 277&-5&11\\
 815&5&-11
 \end{array}
\]
because $5\cdot71+11\cdot67=1092$.  Uniform unit rows have inner product $2/6=1/3$.
Thus
\[
 \frac{\|u_p+u_r\|^2}{\|u_p\|^2+\|u_r\|^2}=\frac43,
 \qquad
 \frac{\|u_p-u_r\|^2}{\|u_p\|^2+\|u_r\|^2}=\frac23. \tag{4}
\]

Equation (4) gives both directions with the same geometry.  Unit normalization removes
the raw mass ratio, but symmetric coefficients amplify and antisymmetric coefficients
save.  Therefore the operator upper bound cannot be reinterpreted as a strict estimate
$G\le(1-\delta)I$ for any positive $\delta$.

The ambient sharpness example in Theorem~\ref{thm:bessel} is intentionally not claimed
to be a literal prime clock.  The Q39 block is the source-model witness needed for the
scoped no-saving statement; it establishes ratio above one without claiming that the
global constant two is attained arithmetically.
